\section{Performance- und Konformitätsanalyse}
\label{sec:performance}

% Umfang: ca. 5 Seiten
% Assignment-Fragen 2-6:
%   2 How long do the different process instances take? Notable patterns?
%   3 Are there any bottlenecks? Where? Reasons?
%   4 How many applications get rejected? Can you find any reasons?
%   5 Can you find any unexpected behavior, e.g. missing permits?
%   6 Are there any patterns that would suggest non-conformance?
% Quellen: Conformance Checking \parencite{carmona_conformance_2018}

\todo{Kapitel schreiben}

\begin{table}[htbp]
\centering
\caption{KPIs Domestic vs International}
\label{tab:kpi_comparison}
\begin{tabular*}{\textwidth}{@{\extracolsep{\fill}} l c c @{}}
\hline
\textbf{KPI} & \textbf{Domestic} & \textbf{International} \\ \hline
Cases & 10.500 & 6.449 \\
Happy-Path-Anteil & 43,98 \% (4.618) & 21,23 \% (1.369) \\
Ø Durchlaufzeit Happy Path & 9 d & 66 d (Ø aller Cases: 75 d ohne Extreme) \\
Konforme Cases & 81 \% (8,47k) & 38 \% (2,47k) \\
Verletzungstypen & 19 & 56 \\
Ø TPT konform vs. verletzend & 10,4 vs. 16,3 d (+5,9 d) & 69,3 vs. 97,1 d (+27,8 d) \\
Schritte pro Case konform vs. verletzend & 5,1 vs. 6,4 & 10,2 vs. 11,8 \\
Rejected by EMPLOYEE & 1.212 Cases (12 \%), +9 d & 1.483 (23 \%), +18 d \\
Rejected by ADMINISTRATION & 846 (8 \%), +11 d & 1.287 (20 \%), +15 d \\
Declaration zweimal hintereinander eingereicht & 776 (7 \%), +12 d & 1.225 (19 \%), +20 d \\ \hline
\end{tabular*}
\end{table}

% ==================================================================
\subsection{Durchlaufzeiten}
\label{sec:durchlaufzeiten}

% Notizen:
% Durchlaufzeit-Verteilung. International liegt der Durchschnitt bei 86 Tagen, mit langem
% Tail bis 743 Tage (ca. 360 Cases im letzten Bin). Celonis nennt zusätzlich 75 Tage
% "ohne Extreme". Nennt beide Werte und erklärt den Unterschied. Die Cases nach 2019 sind
% sehr wenige und wahrscheinlich Ausreißer.
%
% Hinweis: 66 d International beginnen beim Permit -> nicht direkt mit 9 d Domestic
% vergleichbar; fairer Zusatzvergleich: Declaration SUBMITTED -> Payment Handled.

\todo{Durchlaufzeiten schreiben, Screenshot Verteilung (figures/celonis/)}

% ==================================================================
\subsection{Bottlenecks}
\label{sec:bottlenecks}

\textbf{Domestic:} Die Wartezeiten sind kurz. Auffällig ist die Zahlungsphase:
\texttt{FINAL\_APPROVED \textrightarrow\ Request Payment} (3 Arbeitstage) und
\texttt{Request Payment \textrightarrow\ Payment Handled} (4 Arbeitstage), beide bei ca. 95\,\%
der Cases. Das ist ein großer Teil der ca. 9 Tage Happy-Path-Dauer.

\textbf{International:} \texttt{Permit FINAL\_APPROVED \textrightarrow\ Start trip}
dauert 39 Arbeitstage (65\,\% der Cases). Das ist vermutlich kein echter
Engpass, sondern die Vorlaufzeit bis zur geplanten Reise. Der echte Hebel
ist \texttt{End trip \textrightarrow\ Declaration SUBMITTED} mit 13 Arbeitstagen bei
71\,\% der Cases, also die späte Einreichung durch den Mitarbeiter.

\texttt{Permit FINAL\_APPROVED \textrightarrow\ Declaration SUBMITTED}
(17 Arbeitstage, 12\,\% der Cases) betrifft Fälle ohne Reise-Events.
Das sollte kurz erklärt oder nachgeprüft werden.

% ==================================================================
\subsection{Ablehnungen}
\label{sec:rejections}

% Notizen:
% "Rejected by EMPLOYEE" ist plausibel eine Folgeaktivität. In 20 % der International-Cases
% folgt es innerhalb von 3 Arbeitstagen direkt auf "Rejected by ADMINISTRATION". Das stützt
% die Annahme, dass der Mitarbeiter die Rückgabe bestätigt, aber es bleibt eine Hypothese.
%
% Doppelte Einreichungen (19 % bzw. 7 %) sind eine Rework-Schleife, wahrscheinlich nach
% einer Ablehnung.
%
% Offen: Anteil nie bezahlter Cases (kein "Payment Handled"); Gründe stehen nicht im Log
% -> nur Indizien (Betrag, Jahr, Budget-Owner-Beteiligung) per Celonis-Filter vergleichen.

\todo{Ablehnungen je Rolle, Rework, nie bezahlte Cases, mögliche Gründe}

% ==================================================================
\subsection{Unerwartetes Verhalten}
\label{sec:unerwartet}

748 International-Cases (12\,\%) starten mit \texttt{Start trip}, haben
also keinen vorherigen Permit-Antrag, und dauern 29 Tage länger. Das ist
die quantifizierte Antwort auf die Frage nach fehlenden Permits.

637 Cases (10\,\%) haben \texttt{Permit FINAL\_APPROVED by DIRECTOR}
bzw. \texttt{Permit APPROVED by SUPERVISOR}, mit +52 Tagen. Vermutlich
handelt es sich um eine Eskalation bei teuren Reisen; das ist allerdings
eine Hypothese.

% Offen (Jesco, Python): Regel-Checks gegenprüfen -- Trip/Deklaration ohne Permit,
% Doppelzahlungen (Payment Handled > 1), Declaration SUBMITTED vor End trip.

% ==================================================================
\subsection{Muster mangelnder Konformität}
\label{sec:nonconformance}

% Notizen:
% Das Referenzmodell ist nicht normativ. Celonis leitet den Happy Path algorithmisch ab.
% "Undesired activity" heißt nur, dass die Aktivität nicht im Happy Path liegt.
% BUDGET OWNER (27-28 %) oder PRE_APPROVER (6 %) sind deshalb vermutlich legitime Schritte.
% Die 38 % bzw. 81 % konformer Cases solltet ihr als "Abweichung vom häufigsten Pfad"
% formulieren, nicht als Regelverstoß.
%
% Sprung Anfang 2018. In beiden Conformance-Verläufen steigt die Konformität ab 2018-01
% deutlich (Domestic ca. 60 % -> 85 %, International ca. 30 % -> 45 %).
% ACHTUNG: Laut BPIC-2020-Beschreibung umfassen die Daten 2017 nur zwei Fakultäten,
% 2018 die gesamte Universität (Case-Starts Domestic: 2.240 in 2017 vs. 8.260 in 2018).
% Der Anstieg ist daher sehr wahrscheinlich ein Effekt der veränderten Grundgesamtheit,
% nicht (nur) einer Prozessänderung \parencite{van_dongen_bpi_2020}.
%
% Offen (Jesco): Alignments/Fitness gegen normatives Modell -> echte Non-Conformance.

\todo{Non-Conformance-Muster zusammenführen; Alignment-Ergebnisse von Jesco einarbeiten}
